% !TEX root = main.tex
\chapter{Ramsey Degrees and Converse Results}
\label{chap:ramsey-degrees}

In the previous chapters we studied classes with the Ramsey property and saw how ordered Ramsey expansions determine the universal minimal flow of the automorphism group of the Fraïssé limit. However, many natural Fraïssé classes do not themselves satisfy the Ramsey property. The purpose of this chapter is to explain that, even in that situation, there is still a finite combinatorial invariant measuring the extent to which the Ramsey property fails.

This invariant is the \emph{Ramsey degree} of a finite structure. Roughly speaking, if \(\bfa_0\in\cK_0\), then the Ramsey degree of \(\bfa_0\) is the least number \(t\) such that, after passing to a sufficiently large copy of any ambient structure, every colouring of copies of \(\bfa_0\) can be reduced on some substructure to at most \(t\) colours. When \(\cK_0\) admits a reasonable ordered expansion with the Ramsey and ordering properties, these degrees can be computed explicitly in terms of the admissible orderings on \(\bfa_0\).

The chapter has two aims. First, we define these degrees and show how they arise from ordered Ramsey expansions. Second, we prove converse results showing that the combinatorial bounds given by Ramsey degrees recover both the Ramsey property and, ultimately, the universal minimal flow description from Chapter~\ref{chap:examples-umf}.

\section{Patterns and Ramsey degrees}

Throughout this chapter, \(L_0\) is a signature, \(\cK_0\) is a Fraïssé class in \(L_0\), and \(\cK\) is a hereditary order class in the expanded signature \(L=L_0\cup\{<\}\) such that \(\cK\!\restriction\!L_0=\cK_0\).
Thus \(\cK\) is an order expansion of \(\cK_0\).

Let \(\bfa_0\in\cK_0\). We write
\[
    \Adm_{\cK}(\bfa_0)
    =
    \{\prec \mid \prec \text{ is a linear ordering on }A_0 \text{ and } \langle \bfa_0,\prec\rangle\in\cK\}.
\]
So \(\Adm_{\cK}(\bfa_0)\) is the set of all \(\cK\)-admissible linear orderings on the underlying set of \(\bfa_0\).

The group \(\Aut(\bfa_0)\) acts on \(\Adm_{\cK}(\bfa_0)\) in the natural way: if \(g\in\Aut(\bfa_0)\) and \(\prec\in\Adm_{\cK}(\bfa_0)\), then \(g\cdot\prec\) is the linear order on \(A_0\) defined by
\[
    x\,(g\cdot\prec)\,y
    \quad\Longleftrightarrow\quad
    g^{-1}(x)\prec g^{-1}(y).
\]
Since \(\cK\) is closed under isomorphism, we have \(g\cdot\prec\in\Adm_{\cK}(\bfa_0)\) whenever \(\prec\in\Adm_{\cK}(\bfa_0)\).

\begin{definition}
    Let \(\prec\in\Adm_{\cK}(\bfa_0)\). The \emph{\(\bfa_0\)-pattern} of \(\prec\) (with respect to \(\cK\)) is its orbit under the action of \(\Aut(\bfa_0)\) on \(\Adm_{\cK}(\bfa_0)\).
\end{definition}

Equivalently, two orderings \(\prec,\prec'\in\Adm_{\cK}(\bfa_0)\) have the same pattern if and only if \(\langle \bfa_0,\prec\rangle \cong \langle \bfa_0,\prec'\rangle\).
Thus patterns record the isomorphism types of admissible ordered expansions of \(\bfa_0\).
For a fixed \(\bfa_0\), we write
\(\pat_{\cK,\bfa_0}(\prec)\) for the pattern of \(\prec\).

We write
\[
    t_\cK(\bfa_0)
    =
    \left|\Adm_{\cK}(\bfa_0)/\Aut(\bfa_0)\right|
\]
for the number of \(\bfa_0\)-patterns. Since the action of \(\Aut(\bfa_0)\) on \(\Adm_{\cK}(\bfa_0)\) is free, each orbit has cardinality \(|\Aut(\bfa_0)|\). Hence
\[
    t_\cK(\bfa_0)
    =
    \frac{\left|\Adm_{\cK}(\bfa_0)\right|}{\left|\Aut(\bfa_0)\right|}.
\]

To state the relevant partition relation, let
\(\bfa_0\hookrightarrow\bfb_0\hookrightarrow\bfc_0\) in \(\cK_0\), and let
\(k,t\geq 1\). We write
\[
    \bfc_0\to(\bfb_0)^{\bfa_0}_{k,t}
\]
to mean that for every colouring
\[
    c\colon\begin{pmatrix}\bfc_0\\ \bfa_0\end{pmatrix}\to\{1,\dots,k\},
\]
there is some
\[
    \bfb_0'\in\begin{pmatrix}\bfc_0\\ \bfb_0\end{pmatrix}
\]
such that the restriction of \(c\) to
\[
    \begin{pmatrix}\bfb_0'\\ \bfa_0\end{pmatrix}
\]
takes at most \(t\) values. In particular,
\[
    \bfc_0\to(\bfb_0)^{\bfa_0}_{k,1}
    \quad\Longleftrightarrow\quad
    \bfc_0\to(\bfb_0)^{\bfa_0}_{k}.
\]

For later use, suppose \(\bfa_0\hookrightarrow\bfb_0\) in \(\cK_0\) and fix an
ordered expansion \(\langle \bfb_0,\prec^{\bfb_0}\rangle\in\cK\). We write
\[
    t_\cK(\bfa_0,\bfb_0,\prec^{\bfb_0})
\]
for the number of \(\bfa_0\)-patterns represented by orderings
\(\prec\in\Adm_{\cK}(\bfa_0)\) such that
\(\langle \bfa_0,\prec\rangle\hookrightarrow
\langle \bfb_0,\prec^{\bfb_0}\rangle\).
Thus this number counts the patterns of copies of \(\bfa_0\) that actually occur
inside the fixed ordered expansion, and it is at most \(t_\cK(\bfa_0)\).

We can now define the main combinatorial invariant of the chapter.

\begin{definition}
    Let \(\bfa_0\in\cK_0\). The \emph{Ramsey degree} of \(\bfa_0\) in \(\cK_0\),
    denoted \(t(\bfa_0,\cK_0)\), is the least integer \(t\), if it exists, such
    that for every \(\bfb_0\in\cK_0\) with
    \(\bfa_0\hookrightarrow\bfb_0\) and every integer \(k\geq 2\), there is
    \(\bfc_0\in\cK_0\) with \(\bfb_0\hookrightarrow\bfc_0\) and
    \[
        \bfc_0\to(\bfb_0)^{\bfa_0}_{k,t}.
    \]
    If no such integer exists, we put \(t(\bfa_0,\cK_0)=\infty\).
\end{definition}

Thus \(t(\bfa_0,\cK_0)\) measures how far \(\bfa_0\) is from satisfying the usual Ramsey property inside the class \(\cK_0\). The case \(t(\bfa_0,\cK_0)=1\) is exactly the ordinary Ramsey statement for copies of \(\bfa_0\).

The guiding principle of this chapter is that, when \(\cK_0\) admits a suitable ordered Ramsey expansion \(\cK\), the Ramsey degree of \(\bfa_0\) is determined by the finite set of admissible patterns on \(\bfa_0\). In the next section we make this precise.

A basic example is given by the class \(\mathcal{GR}\) of finite graphs. This class does not have the Ramsey property, but its standard ordered expansion by linear orders does. If \(\bfa_0\in\mathcal{GR}\), then every linear ordering of \(A_0\) is admissible, and two such orderings have the same pattern exactly when they differ by an automorphism of \(\bfa_0\). Thus
\[
    t(\bfa_0,\mathcal{GR})
    =
    \frac{|A_0|!}{|\Aut(\bfa_0)|}.
\]
This illustrates the general picture: Ramsey degrees arise from counting the admissible order-types of a structure up to automorphism.

\section{Ramsey degrees from Ramsey expansions}

We now show that if the ordered class \(\cK\) has the Ramsey property, then the number of admissible patterns gives an upper bound on the number of colours that must remain on a large substructure. If \(\cK\) also has the ordering property, then this bound is optimal. In particular, the Ramsey degree of \(\bfa_0\) is exactly the number of \(\cK\)-patterns on \(\bfa_0\).

\begin{proposition}
    \label{prop:ramsey-expansion-upper-bound}
    Assume that \(\cK\) has the Ramsey property. Let
    \(\bfa_0\hookrightarrow\bfb_0\) in \(\cK_0\), let
    \(\langle \bfb_0,\prec^{\bfb_0}\rangle\in\cK\), and let \(k\geq 2\). Then there
    is \(\bfc_0\) in \(\cK_0\) with \(\bfb_0\hookrightarrow\bfc_0\) such that
    \[
        \bfc_0\to(\bfb_0)^{\bfa_0}_{k,\,t_\cK(\bfa_0,\bfb_0,\prec^{\bfb_0})}.
    \]
    In particular, for every hereditary class \(\cK'\subseteq\cK\) with
    \(\cK'\!\restriction\!L_0=\cK_0\), there is \(\bfc_0\) in \(\cK_0\) with
    \(\bfb_0\hookrightarrow\bfc_0\) such that
    \[
        \bfc_0\to(\bfb_0)^{\bfa_0}_{k,\,t_{\cK'}(\bfa_0)}.
    \]
\end{proposition}

\begin{proof}
    Put \(t=t_\cK(\bfa_0,\bfb_0,\prec^{\bfb_0})\). Choose representatives \(\prec_1,\dots,\prec_t\) for the \(\bfa_0\)-patterns realised inside \(\langle\bfb_0,\prec^{\bfb_0}\rangle\).

    We define inductively ordered structures
    \[
        \langle\bfd_0^{(i)},\prec^{\bfd_0^{(i)}}\rangle\in\cK
        \qquad (0\leq i\leq t)
    \]
    as follows. Let
    \(\langle\bfd_0^{(0)},\prec^{\bfd_0^{(0)}}\rangle
      =\langle\bfb_0,\prec^{\bfb_0}\rangle\), and, for each
    \(1\leq i\leq t\), choose
    \(\langle\bfd_0^{(i)},\prec^{\bfd_0^{(i)}}\rangle\in\cK\)
    such that
    \[
        \langle\bfd_0^{(i)},\prec^{\bfd_0^{(i)}}\rangle\to
        \bigl(\langle\bfd_0^{(i-1)},\prec^{\bfd_0^{(i-1)}}\rangle\bigr)^{\langle\bfa_0,\prec_i\rangle}_{k}.
    \]
    Set \(\bfc_0=\bfd_0^{(t)}\) and write
    \(\prec^{\bfc_0}=\prec^{\bfd_0^{(t)}}\).

    We claim that \(\bfc_0\) has the required property. Let
    \[
        c\colon\begin{pmatrix}\bfc_0\\ \bfa_0\end{pmatrix}\to\{1,\dots,k\}
    \]
    be any colouring of copies of \(\bfa_0\) in \(\bfc_0\).

    We now work downward from
    \(\langle\bfd_0^{(t)},\prec^{\bfd_0^{(t)}}\rangle\). For the pattern
    represented by \(\prec_t\), define a colouring
    \[
        c_t\colon
        \begin{pmatrix}
            \langle\bfd_0^{(t)},\prec^{\bfd_0^{(t)}}\rangle \\
            \langle\bfa_0,\prec_t\rangle
        \end{pmatrix}
        \to
        \{1,\dots,k\}
    \]
    by sending a copy
    \[
        \langle\bfa_0',\prec^{\bfd_0^{(t)}}\restriction A_0'\rangle\cong\langle\bfa_0,\prec_t\rangle
    \]
    to the colour \(c(\bfa_0')\) of its reduct. By the choice of
    \(\langle\bfd_0^{(t)},\prec^{\bfd_0^{(t)}}\rangle\), there is a copy
    \[
        \langle\widetilde{\bfd}_0^{(t-1)},
        \prec^{\widetilde{\bfd}_0^{(t-1)}}\rangle\cong
        \langle\bfd_0^{(t-1)},\prec^{\bfd_0^{(t-1)}}\rangle
    \]
    inside \(\langle\bfd_0^{(t)},\prec^{\bfd_0^{(t)}}\rangle\) on which every
    copy of \(\langle\bfa_0,\prec_t\rangle\) has the same \(c_t\)-colour,
    say \(k_t\).

    Next define
    \[
        c_{t-1}\colon
        \begin{pmatrix}
            \langle\widetilde{\bfd}_0^{(t-1)},
            \prec^{\widetilde{\bfd}_0^{(t-1)}}\rangle \\
            \langle\bfa_0,\prec_{t-1}\rangle
        \end{pmatrix}
        \to
        \{1,\dots,k\}
    \]
    in the same way, by pulling back the colouring \(c\) from reducts. Again by the Ramsey property, there is a copy
    \[
        \langle\widetilde{\bfd}_0^{(t-2)},
        \prec^{\widetilde{\bfd}_0^{(t-2)}}\rangle\cong
        \langle\bfd_0^{(t-2)},\prec^{\bfd_0^{(t-2)}}\rangle
    \]
    inside \(\langle\widetilde{\bfd}_0^{(t-1)},
    \prec^{\widetilde{\bfd}_0^{(t-1)}}\rangle\) on which every copy of
    \(\langle\bfa_0,\prec_{t-1}\rangle\) has the same colour, say \(k_{t-1}\).

    Continuing in this way, after \(t\) steps we obtain a copy
    \[
        \langle\bfb_0',\prec^{\bfb_0'}\rangle\cong
        \langle\bfb_0,\prec^{\bfb_0}\rangle
    \]
    inside \(\langle\bfc_0,\prec^{\bfc_0}\rangle\) together with colours
    \(k_1,\dots,k_t\in\{1,\dots,k\}\) such that, for each
    \(i\in\{1,\dots,t\}\), every copy of \(\langle\bfa_0,\prec_i\rangle\)
    inside \(\langle\bfb_0',\prec^{\bfb_0'}\rangle\) has colour \(k_i\).

    Now let
    \[
        \bfa_0'\in\begin{pmatrix}\bfb_0'\\ \bfa_0\end{pmatrix}.
    \]
    The induced ordering \(\prec^{\bfb_0'}\restriction A_0'\) realises one of the patterns represented by some \(\prec_i\). Hence \(c(\bfa_0')=k_i\). Therefore the colouring \(c\) takes at most the \(t\) values \(k_1,\dots,k_t\)
    on
    \[
        \begin{pmatrix}\bfb_0'\\ \bfa_0\end{pmatrix}.
    \]
    This proves
    \[
        \bfc_0\to(\bfb_0)^{\bfa_0}_{k,t}.
    \]

    For the final claim, note that \(t_\cK(\bfa_0,\bfb_0,\prec^{\bfb_0})\leq t_{\cK'}(\bfa_0)\),
    since every pattern realised in \(\langle\bfb_0,\prec^{\bfb_0}\rangle\) with respect to \(\cK'\) is also a \(\cK'\)-pattern on \(\bfa_0\). Hence the same \(\bfc_0\) also satisfies
    \[
        \bfc_0\to(\bfb_0)^{\bfa_0}_{k,\,t_{\cK'}(\bfa_0)}.
    \]
\end{proof}

The proof yields a slightly stronger conclusion: one may choose a linear ordering on \(\bfc_0\) such that, on a suitable copy of \(\bfb_0\), the colour of a copy of \(\bfa_0\) depends only on the pattern induced by that ordering. We will not isolate this as a separate proposition, since the weaker statement above is all that we need.

\begin{proposition}
    \label{prop:ordering-property-lower-bound}
    Assume that \(\cK\) has both the Ramsey property and the ordering property.
    Then for every \(\bfa_0\in\cK_0\) there exists \(\bfb_0\in\cK_0\) with
    \(\bfa_0\hookrightarrow\bfb_0\) such that for every \(\bfc_0\in\cK_0\)
    with \(\bfb_0\hookrightarrow\bfc_0\) there is a colouring
    \[
        c\colon\begin{pmatrix}\bfc_0\\ \bfa_0\end{pmatrix}\to\{1,\dots,t_\cK(\bfa_0)\}
    \]
    with the following property: for every
    \[
        \bfb_0'\in\begin{pmatrix}\bfc_0\\ \bfb_0\end{pmatrix},
    \]
    the colouring \(c\) takes all the values \(1,\dots,t_\cK(\bfa_0)\) on
    \[
        \begin{pmatrix}\bfb_0'\\ \bfa_0\end{pmatrix}.
    \]
\end{proposition}

\begin{proof}
    By the ordering property, there is \(\bfb_0\in\cK_0\) with
    \(\bfa_0\hookrightarrow\bfb_0\) such that whenever
    \[
        \langle\bfa_0,\prec\rangle,\langle\bfb_0,\prec'\rangle\in\cK,
    \]
    we have
    \[
        \langle\bfa_0,\prec\rangle\hookrightarrow\langle\bfb_0,\prec'\rangle.
    \]
    Equivalently, every admissible \(\bfa_0\)-pattern is realised by some copy of \(\bfa_0\) inside every ordered expansion of \(\bfb_0\) belonging to \(\cK\).

    Now let \(\bfc_0\in\cK_0\) with \(\bfb_0\hookrightarrow\bfc_0\), and choose a linear ordering \(\prec^{\bfc_0}\) on \(C_0\) such that \(\langle\bfc_0,\prec^{\bfc_0}\rangle\in\cK\). Enumerate the \(\bfa_0\)-patterns as \(p_1,\dots,p_{t_\cK(\bfa_0)}\).
    Define
    \[
        c\colon\begin{pmatrix}\bfc_0\\ \bfa_0\end{pmatrix}\to\{1,\dots,t_\cK(\bfa_0)\}
    \]
    by declaring that, for
    \[
        \bfa_0'\in\begin{pmatrix}\bfc_0\\ \bfa_0\end{pmatrix},
    \]
    we put \(c(\bfa_0')=i\) if the ordered structure \(\langle\bfa_0',\prec^{\bfc_0}\restriction A_0'\rangle\) has pattern \(p_i\).

    Let
    \[
        \bfb_0'\in\begin{pmatrix}\bfc_0\\ \bfb_0\end{pmatrix}.
    \]
    Then \(\langle\bfb_0',\prec^{\bfc_0}\restriction B_0'\rangle\in\cK\), and by the choice of \(\bfb_0'\), every admissible \(\bfa_0\)-pattern is realised inside this ordered structure. Therefore the colouring \(c\) takes every value \(1,\dots,t_\cK(\bfa_0)\)
    on
    \[
        \begin{pmatrix}\bfb_0'\\ \bfa_0\end{pmatrix}.
    \]
\end{proof}

\begin{corollary}
    \label{cor:ramsey-degree-pattern-count}
    Assume that \(\cK\) has the Ramsey property and the ordering property. Then for every \(\bfa_0\in\cK_0\),
    \[
        t(\bfa_0,\cK_0)=t_\cK(\bfa_0).
    \]
    Equivalently, \(t_\cK(\bfa_0)\) is the least integer \(t\) such that for every \(\bfb_0\in\cK_0\) with \(\bfa_0\hookrightarrow\bfb_0\) and every \(k\geq 2\), there exists \(\bfc_0\in\cK_0\) with \(\bfb_0\hookrightarrow\bfc_0\) and
    \[
        \bfc_0\to(\bfb_0)^{\bfa_0}_{k,t}.
    \]
\end{corollary}

\begin{proof}
    Proposition~\ref{prop:ramsey-expansion-upper-bound} shows that \(t_\cK(\bfa_0)\) has the required partition property, so
    \[
        t(\bfa_0,\cK_0)\leq t_\cK(\bfa_0).
    \]
    On the other hand, Proposition~\ref{prop:ordering-property-lower-bound} shows that no smaller number can work. Hence
    \[
        t(\bfa_0,\cK_0)=t_\cK(\bfa_0).
    \]
\end{proof}

In particular, whenever \(\cK\) has the Ramsey and ordering properties, the number of admissible patterns on \(\bfa_0\) depends only on the reduct class \(\cK_0\), not on the particular expansion \(\cK\). Thus ordered Ramsey expansions do not merely imply the finiteness of Ramsey degrees; they compute those degrees exactly.

\section{Converse criteria}

We now prove converses to the results of the previous section. The first shows that, under the ordering property, the partition relation with bound \(t_\cK(\bfa_0)\) already forces the Ramsey property. We then isolate a local form of the ordering property and show that, assuming the Ramsey property, it is equivalent to the equality \(t(\bfa_0,\cK_0)=t_\cK(\bfa_0)\).

\begin{proposition}
    \label{prop:partition-bound-implies-ramsey}
    Suppose that \(\cK\) is reasonable and has the ordering property. Assume
    moreover that for every \(\bfa_0\hookrightarrow\bfb_0\) in \(\cK_0\) and
    every \(k\geq 2\), there exists \(\bfc_0\in\cK_0\) with
    \(\bfb_0\hookrightarrow\bfc_0\) such that
    \[
        \bfc_0\to(\bfb_0)^{\bfa_0}_{k,\,t_\cK(\bfa_0)}.
    \]
    Then \(\cK\) has the Ramsey property.

    In particular, any reasonable order class \(\cK\) with the ordering property that is contained in an expansion of \(\cK_0\) with the Ramsey property also has the Ramsey property.
\end{proposition}

\begin{proof}
    Let
    \[
        \langle\bfa_0,\prec\rangle,\langle\bfb_0,\prec'\rangle\in\cK
        \qquad\text{with}\qquad
        \langle\bfa_0,\prec\rangle\hookrightarrow\langle\bfb_0,\prec'\rangle.
    \]
    We must find \(\langle\bfc_0,\prec''\rangle\in\cK\) such that
    \[
        \langle\bfc_0,\prec''\rangle\to
        \bigl(\langle\bfb_0,\prec'\rangle\bigr)^{\langle\bfa_0,\prec\rangle}_{2}.
    \]

    By the ordering property, there exists \(\bfb_1\in\cK_0\) with
    \(\bfb_0\hookrightarrow\bfb_1\) such that whenever
    \[
        \langle\bfb_0,\prec^{\bfb_0}\rangle,\langle\bfb_1,\prec^{\bfb_1}\rangle\in\cK,
    \]
    we have
    \[
        \langle\bfb_0,\prec^{\bfb_0}\rangle\hookrightarrow\langle\bfb_1,\prec^{\bfb_1}\rangle.
    \]
    Since \(\cK\) is reasonable and \(\langle\bfa_0,\prec\rangle\hookrightarrow\langle\bfb_0,\prec'\rangle\), it follows that for every ordering \(\prec^{\bfb_1}\) on \(B_1\) with \(\langle\bfb_1,\prec^{\bfb_1}\rangle\in\cK\), we also have
    \[
        \langle\bfa_0,\prec\rangle\hookrightarrow\langle\bfb_1,\prec^{\bfb_1}\rangle.
    \]

    Now apply the hypothesis with \(\bfa_0\), \(\bfb_1\), and \(k=t_\cK(\bfa_0)+1\).
    Thus there exists \(\bfc_0\in\cK_0\) with
    \(\bfb_1\hookrightarrow\bfc_0\) such that
    \[
        \bfc_0\to(\bfb_1)^{\bfa_0}_{\,t_\cK(\bfa_0)+1,\;t_\cK(\bfa_0)}.
    \]
    Choose an ordering \(\prec''\) on \(C_0\) such that \(\langle\bfc_0,\prec''\rangle\in\cK\).
    We claim that this ordered structure has the required Ramsey property.

    Let
    \[
        c\colon
        \begin{pmatrix}
            \langle\bfc_0,\prec''\rangle \\
            \langle\bfa_0,\prec\rangle
        \end{pmatrix}
        \to
        \{1,2\}
    \]
    be a colouring of copies of \(\langle\bfa_0,\prec\rangle\) in \(\langle\bfc_0,\prec''\rangle\). Enumerate the \(\bfa_0\)-patterns as \(p_1,\dots,p_{t_\cK(\bfa_0)}\).
    Using \(c\), define a colouring
    \[
        c'\colon
        \begin{pmatrix}\bfc_0\\ \bfa_0\end{pmatrix}
        \to
        \{0,1,\dots,t_\cK(\bfa_0)\}
    \]
    as follows. For
    \[
        \bfa_0'\in\begin{pmatrix}\bfc_0\\ \bfa_0\end{pmatrix},
    \]
    let
    \[
        q(\bfa_0')
        =
        \pat_{\cK,\bfa_0}(\prec''\restriction A_0')
    \]
    be the pattern of the induced ordering on \(A_0'\). Set
    \[
        c'(\bfa_0')=
        \begin{cases}
            0, & \text{if }q(\bfa_0')=\pat_{\cK,\bfa_0}(\prec)
            \text{ and }c(\langle\bfa_0',\prec''\restriction A_0'\rangle)=2,  \\[1ex]
            i, & \text{if }q(\bfa_0')=p_i\text{ and either }
            p_i\neq\pat_{\cK,\bfa_0}(\prec)
            \text{ or }c(\langle\bfa_0',\prec''\restriction A_0'\rangle)=1.
        \end{cases}
    \]
    This uses at most \(t_\cK(\bfa_0)+1\) colours.

    Hence there exists
    \[
        \bfb_1'\in\begin{pmatrix}\bfc_0\\ \bfb_1\end{pmatrix}
    \]
    such that \(c'\) takes at most \(t_\cK(\bfa_0)\) values on
    \[
        \begin{pmatrix}\bfb_1'\\ \bfa_0\end{pmatrix}.
    \]
    Since \(\langle\bfb_1',\prec''\restriction B_1'\rangle\in\cK\) and \(\langle\bfa_0,\prec\rangle\) embeds into every ordered expansion of \(\bfb_1\) in \(\cK\), the pattern of \(\prec\) is realised by some copy of \(\bfa_0\) inside \(\bfb_1'\). Moreover, every \(\bfa_0\)-pattern realised in \(\langle\bfb_1',\prec''\restriction B_1'\rangle\) contributes at least one value of \(c'\) on
    \[
        \begin{pmatrix}\bfb_1'\\ \bfa_0\end{pmatrix}.
    \]
    Since \(c'\) takes at most \(t_\cK(\bfa_0)\) values there, the value \(0\) must occur and, for the pattern of \(\prec\), no additional value can occur. It follows that all copies of \(\langle\bfa_0,\prec\rangle\) inside \(\langle\bfb_1',\prec''\restriction B_1'\rangle\) have the same \(c\)-colour, namely \(2\).

    Finally, by the choice of \(\bfb_1\) we have
    \[
        \langle\bfb_0,\prec'\rangle\hookrightarrow\langle\bfb_1',\prec''\restriction B_1'\rangle,
    \]
    so
    \[
        \langle\bfc_0,\prec''\rangle\to
        \bigl(\langle\bfb_0,\prec'\rangle\bigr)^{\langle\bfa_0,\prec\rangle}_{2}.
    \]
    Thus \(\cK\) has the Ramsey property.

    For the final assertion, if \(\cK\) is contained in an expansion of \(\cK_0\) with the Ramsey property, then Proposition~\ref{prop:ramsey-expansion-upper-bound} yields exactly the partition relation assumed above. Hence the first part applies.
\end{proof}

To formulate the converse to Corollary~\ref{cor:ramsey-degree-pattern-count}, it is useful to isolate a local form of the ordering property.

\begin{definition}
    Let \(\bfa_0\in\cK_0\). We say that \(\cK\) has the \emph{ordering property at} \(\bfa_0\) if there exists \(\bfb_0\in\cK_0\) with \(\bfa_0\hookrightarrow\bfb_0\) such that for every pair of orderings \(\prec\) on \(A_0\) and \(\prec'\) on \(B_0\) with
    \[
        \langle\bfa_0,\prec\rangle,\langle\bfb_0,\prec'\rangle\in\cK,
    \]
    we have
    \[
        \langle\bfa_0,\prec\rangle\hookrightarrow\langle\bfb_0,\prec'\rangle.
    \]
\end{definition}

Thus \(\cK\) has the ordering property if and only if it has the ordering property at every \(\bfa_0\in\cK_0\).

\begin{proposition}
    \label{prop:local-ordering-property-ramsey-degree}
    Assume that \(\cK\) has the Ramsey property. Then, for each \(\bfa_0\in\cK_0\), the following are equivalent:

    \begin{enumerate}
        \item[(i)] \(\cK\) has the ordering property at \(\bfa_0\);
        \item[(ii)] \(t(\bfa_0,\cK_0)=t_\cK(\bfa_0)\).
    \end{enumerate}
\end{proposition}

\begin{proof}
    Assume first that \(\cK\) has the ordering property at \(\bfa_0\). Then the proof of Proposition~\ref{prop:ordering-property-lower-bound} shows that there exists \(\bfb_0\) with \(\bfa_0\hookrightarrow\bfb_0\) such that every admissible \(\bfa_0\)-pattern is realised inside every ordered expansion of \(\bfb_0\) belonging to \(\cK\). Combining this with Proposition~\ref{prop:ramsey-expansion-upper-bound}, exactly as in the proof of Corollary~\ref{cor:ramsey-degree-pattern-count}, we obtain \(t(\bfa_0,\cK_0)=t_\cK(\bfa_0)\).

    Conversely, suppose that \(t(\bfa_0,\cK_0)=t_\cK(\bfa_0)\)
    but that \(\cK\) does not have the ordering property at \(\bfa_0\). Then for every \(\bfb_0\in\cK_0\) with \(\bfa_0\hookrightarrow\bfb_0\) there exists an ordering \(\prec'\) on \(B_0\) such that \(\langle\bfb_0,\prec'\rangle\in\cK\) and some admissible ordering \(\prec\) on \(A_0\) with \(\langle\bfa_0,\prec\rangle\in\cK\)
    fails to embed into \(\langle\bfb_0,\prec'\rangle\). Hence the number of \(\bfa_0\)-patterns realised inside \(\langle\bfb_0,\prec'\rangle\) satisfies
    \[
        t_\cK(\bfa_0,\bfb_0,\prec')\leq t_\cK(\bfa_0)-1.
    \]
    Applying Proposition~\ref{prop:ramsey-expansion-upper-bound} with this ordered structure, we obtain, for every \(k\geq 2\), some \(\bfc_0\) with \(\bfb_0\hookrightarrow\bfc_0\) such that
    \[
        \bfc_0\to(\bfb_0)^{\bfa_0}_{k,\,t_\cK(\bfa_0)-1}.
    \]
    This implies \(t(\bfa_0,\cK_0)\leq t_\cK(\bfa_0)-1\),
    contrary to the assumption. Therefore \(\cK\) has the ordering property at \(\bfa_0\).
\end{proof}

\begin{corollary}
    \label{cor:ordering-property-ramsey-degree}
    Assume that \(\cK\) has the Ramsey property. Then the following are equivalent:

    \begin{enumerate}
        \item[(i)] \(\cK\) has the ordering property;
        \item[(ii)] for every \(\bfa_0\in\cK_0\),
              \[
                  t(\bfa_0,\cK_0)=t_\cK(\bfa_0).
              \]
    \end{enumerate}
\end{corollary}

\begin{proof}
    This follows immediately from Proposition~\ref{prop:local-ordering-property-ramsey-degree}, since the ordering property is equivalent to having the ordering property at each \(\bfa_0\in\cK_0\).
\end{proof}


\section{Recovering Ramsey expansions}

We now return to the dynamical viewpoint of the previous chapters. The results of this section show that the combinatorial information encoded by Ramsey degrees is strong enough to recover ordered expansions with the Ramsey and ordering properties, and hence to recover the universal minimal flow description from Chapter~\ref{chap:examples-umf}.

\begin{theorem}
    \label{thm:recover-ramsey-ordering-subexpansion}
    Let \(\cK_0\) be a Fraïssé class in a signature \(L_0\), and let \(\cK\) be a reasonable Fraïssé order class in the expanded signature \(L=L_0\cup\{<\}\) such that \(\cK\!\restriction\!L_0=\cK_0\). Assume that \(\cK\) has the Ramsey property. Then there exists a reasonable Fraïssé order class \(\cK'\subseteq\cK\) such that \(\cK'\!\restriction\!L_0=\cK_0\) and \(\cK'\) has both the Ramsey property and the ordering property.
\end{theorem}

\begin{proof}
    Let \(\bff_0=\Flim(\cK_0)\) and \(G_0=\Aut(\bff_0)\). Consider the \(G_0\)-flow \(X_\cK\) of all \(\cK\)-admissible linear orderings on \(F_0\). Choose a minimal subflow \(X'\subseteq X_\cK\), and fix some \(\prec'\in X'\). Let \(\cK'=\Age(\langle \bff_0,\prec'\rangle)\).
    Then \(\cK'\subseteq\cK\), and since every finite \(L_0\)-substructure of \(\bff_0\) lies in \(\cK_0\), we have \(\cK'\!\restriction\!L_0=\cK_0\).
    Moreover, \(\cK'\) is hereditary and has the joint embedding property, and it is easy to verify that \(\cK'\) is reasonable.

    We claim that
    \[
        X'=X_{\cK'}.
    \]
    Indeed, every finite substructure of \(\langle\bff_0,\prec'\rangle\) belongs to \(\cK'\), so every ordering in the orbit closure \(\overline{G_0\cdot \prec'}=X'\) is \(\cK'\)-admissible. Conversely, if an ordering on \(F_0\) is \(\cK'\)-admissible, then every finite restriction occurs in \(\Age(\langle\bff_0,\prec'\rangle)\), hence belongs to the orbit closure of \(\prec'\). Thus the two spaces coincide.

    Since \(X'=X_{\cK'}\) is minimal, it follows from
    Theorem~\ref{thm:minimality-ordering-property} that \(\cK'\) has the ordering
    property.

    It remains to show that \(\cK'\) has the Ramsey property. By Proposition~\ref{prop:partition-bound-implies-ramsey}, it is enough to verify that for every \(\bfa_0\hookrightarrow\bfb_0\) in \(\cK_0\) and every \(k\geq 2\), there exists \(\bfc_0\in\cK_0\) with \(\bfb_0\hookrightarrow\bfc_0\) such that
    \[
        \bfc_0\to(\bfb_0)^{\bfa_0}_{k,\,t_{\cK'}(\bfa_0)}.
    \]
    But this follows from Proposition~\ref{prop:ramsey-expansion-upper-bound}, since \(\cK\) has the Ramsey property and \(\cK'\subseteq\cK\) is a hereditary expansion of \(\cK_0\).

    Finally, because \(\cK'\) is reasonable and has the Ramsey and ordering properties, Corollary~\ref{cor:ordering-property-ramsey-degree} implies that the numbers \(t_{\cK'}(\bfa_0)\) are the Ramsey degrees in \(\cK_0\). In particular, Proposition~\ref{prop:partition-bound-implies-ramsey} yields the Ramsey property for \(\cK'\), and hence \(\cK'\) is a reasonable Fraïssé order class with the required properties.
\end{proof}

The preceding theorem shows that, once one has a reasonable Ramsey expansion, one may pass to a smaller expansion in which minimality is built in. The next result is the converse direction of the universal-minimal-flow theorem proved earlier.

\begin{theorem}
    \label{thm:umf-iff-ramsey-ordering}
    Let \(L\supseteq\{<\}\) be a signature, let \(L_0=L\setminus\{<\}\), and let \(\cK\) be a reasonable Fraïssé order class in \(L\). Put
    \[
        \cK_0=\cK\!\restriction\!L_0,
        \qquad
        \bff=\Flim(\cK),
        \qquad
        \bff_0=\bff\!\restriction\!L_0=\Flim(\cK_0),
    \]
    and let \(G_0=\Aut(\bff_0)\) and \(G=\Aut(\bff)\).
    Let \(X_\cK\) be the compact space of all \(\cK\)-admissible linear orderings on the underlying set \(F_0\) of \(\bff_0\). Then the following are equivalent:

    \begin{enumerate}
        \item[(i)] \(\cK\) has the Ramsey property and the ordering property;
        \item[(ii)] \(X_\cK\) is the universal minimal flow of \(G_0\).
    \end{enumerate}
\end{theorem}

\begin{proof}
    The implication \((i)\Rightarrow(ii)\) is
    Corollary~\ref{cor:ramsey-ordering-umf}.

    Now assume that \(X_\cK\) is the universal minimal flow of \(G_0\). Since it is
    in particular minimal, Theorem~\ref{thm:minimality-ordering-property}
    implies that \(\cK\) has the ordering property. It therefore remains to prove
    that \(\cK\) has the Ramsey property.

    By Proposition~\ref{prop:partition-bound-implies-ramsey}, it is enough to show that for every \(\bfa_0\hookrightarrow\bfb_0\) in \(\cK_0\),
    every \(k\geq 2\), and every colouring
    \[
        c\colon\begin{pmatrix}\bff_0\\ \bfa_0\end{pmatrix}\to\{1,\dots,k\},
    \]
    there exists
    \[
        \bfb_0'\in\begin{pmatrix}\bff_0\\ \bfb_0\end{pmatrix}
    \]
    such that the restriction of \(c\) to
    \[
        \begin{pmatrix}\bfb_0'\\ \bfa_0\end{pmatrix}
    \]
    takes at most \(t_\cK(\bfa_0)\) values.

    Consider the compact \(G_0\)-flow
    \[
        \{1,\dots,k\}^{\begin{pmatrix}\bff_0\\ \bfa_0\end{pmatrix}},
    \]
    where \(G_0\) acts by
    \[
        (g\cdot \gamma)(\bfa_0')=\gamma(g^{-1}(\bfa_0'))
        \qquad
        \left(g\in G_0,\ \gamma\in \{1,\dots,k\}^{\begin{pmatrix}\bff_0\\ \bfa_0\end{pmatrix}}\right).
    \]
    Let \(X=\overline{G_0\cdot c}\) be the orbit closure of the colouring \(c\). Since \(X_\cK\) is universal minimal, there is a \(G_0\)-map \(\pi\colon X_\cK\to X\). Let \(\prec_0\) be the distinguished linear ordering on \(F_0\) coming from the Fraïssé limit \(\bff=\langle\bff_0,\prec_0\rangle\), and put \(c_0=\pi(\prec_0)\).

    Because every element of \(G=\Aut(\bff)\) fixes the ordering \(\prec_0\), and \(\pi\) is \(G_0\)-equivariant, every element of \(G\) also fixes \(c_0\). Hence if
    \[
        \bfa_0',\bfa_0''\in\begin{pmatrix}\bff_0\\ \bfa_0\end{pmatrix}
    \]
    have the same \(\cK\)-pattern with respect to the ordering \(\prec_0\), then some element of \(G\) sends \(\bfa_0'\) to \(\bfa_0''\), and therefore \(c_0(\bfa_0')=c_0(\bfa_0'')\). Thus the value of \(c_0\) on a copy of \(\bfa_0\) depends only on its pattern. It follows that \(c_0\) takes at most \(t_\cK(\bfa_0)\) values.

    Now fix any
    \[
        \bfb_0''\in\begin{pmatrix}\bff_0\\ \bfb_0\end{pmatrix}.
    \]
    Since \(c_0\in \overline{G_0\cdot c}\), there exists \(g\in G_0\) such that
    \[
        (g\cdot c)\restriction\begin{pmatrix}\bfb_0''\\ \bfa_0\end{pmatrix}
        =
        c_0\restriction\begin{pmatrix}\bfb_0''\\ \bfa_0\end{pmatrix}.
    \]
    Let \(\bfb_0'=g^{-1}(\bfb_0'')\).
    Then the restriction of \(c\) to
    \[
        \begin{pmatrix}\bfb_0'\\ \bfa_0\end{pmatrix}
    \]
    has the same set of values as the restriction of \(c_0\) to
    \[
        \begin{pmatrix}\bfb_0''\\ \bfa_0\end{pmatrix},
    \]
    and therefore takes at most \(t_\cK(\bfa_0)\) values. This is exactly the partition relation required in Proposition~\ref{prop:partition-bound-implies-ramsey}.

    Hence \(\cK\) has the Ramsey property. Since we already know that \(\cK\) has the ordering property, this proves \((ii)\Rightarrow(i)\).
\end{proof}

The preceding two theorems complete the converse direction of the Kechris--Pestov--Todorcevic (KPT) correspondence~\cite{KPT05} developed in the previous chapters. Ordered Ramsey expansions not only produce universal minimal flows, but can also be recovered from the dynamical structure of those flows. In this sense, Ramsey degrees measure the residual combinatorial complexity of the reduct class, while the Ramsey and ordering properties together characterise the expansions that realise the universal minimal flow.
